% concatenated from main.tex
\documentclass[12pt]{amsart}
\usepackage{mathtools,amssymb,amsfonts}
\usepackage{amsthm}
\usepackage{hyperref}

\hypersetup{
  colorlinks=true,
  linkcolor=blue,
  citecolor=blue,
  urlcolor=blue,
  pdftitle={Critical-point-free energy for fractional-Toledo representations},
  pdfauthor={Rivu Bardhan, Anu Dhochak, and Pradip Kumar}
}

\newtheorem{theorem}{Theorem}[section]
\newtheorem{proposition}[theorem]{Proposition}
\newtheorem{lemma}[theorem]{Lemma}
\newtheorem{corollary}[theorem]{Corollary}
\theoremstyle{remark}
\newtheorem{remark}[theorem]{Remark}

\newcommand{\CH}{\mathbb{CH}}
\newcommand{\Teich}{\mathcal T}
\newcommand{\cX}{\mathcal X}
\newcommand{\PU}{\mathrm{PU}}
\newcommand{\SU}{\mathrm{SU}}
\newcommand{\tr}{\operatorname{tr}}
\newcommand{\im}{\operatorname{im}}
\newcommand{\Hess}{\operatorname{Hess}}
\newcommand{\II}{\mathrm{II}}
\newcommand{\dd}{\,\mathrm d}

\title[Critical-point-free energy]{Critical-point-free energy for\\
fractional-Toledo representations}

\author[R. Bardhan]{Rivu Bardhan}
\address{Department of Mathematics, Shiv Nadar University, Dadri 201314,
Uttar Pradesh, India}
\email{rb212@snu.edu.in}

\author[A. Dhochak]{Anu Dhochak}
\address{International Centre for Theoretical Sciences, Tata Institute of
Fundamental Research, Bengaluru 560089, India}
\email{anu.dhochak@icts.res.in}

\author[P. Kumar]{Pradip Kumar}
\address{Department of Mathematics, Shiv Nadar University, Dadri 201314,
Uttar Pradesh, India}
\email{pradip.kumar@snu.edu.in}
\date{}

\subjclass[2020]{Primary 53C43, 53C42, 20H10; Secondary 32G15, 58E20}

\begin{document}

\begin{abstract}
Let $S_g$ be a closed oriented surface of genus $g\ge2$.  For a reductive
representation $\rho:\pi_1(S_g)\to\PU(2,1)$, let $E_\rho$ be the energy
function on Teichm\"uller space associated to equivariant harmonic maps into
$\CH^2$.  For every positive integer $d$ with $3\nmid d$, all sufficiently
large $h$, and every $g>h$, we construct an irreducible reductive
representation
\[
\rho_{g,h,d}:\pi_1(S_g)\to\PU(2,1)
\]
with
\[
\tau(\rho_{g,h,d})=2h-2-\frac{2d}{3}\notin\mathbb Z,
\qquad
\operatorname{Crit}(E_{\rho_{g,h,d}})=\varnothing.
\]
Consequently, the associated branched-minimal-surface forgetful map is not
surjective in these nonintegral Toledo components.
\end{abstract}

\maketitle

\section{Introduction}

Let $S_g$ be a closed oriented surface of genus $g\ge2$, and equip
$\CH^2$ with the symmetric metric with sectional curvatures in
$[-4,-1]$. 
If $\rho:\pi_1(S_g)\to\PU(2,1)$ is reductive and
$J\in\Teich(S_g)$, Corlette's theorem
\cite[Corollary~3.5]{Corlette} gives a $\rho$-equivariant harmonic map
\[
 f_{\rho,J}:(\widetilde S_g,J)\longrightarrow\CH^2.
\]
When $\rho$ is irreducible, this map is nonconstant and cannot have image in a
geodesic; it is therefore unique by the equivariant uniqueness theorem
\cite[Theorem~3.2]{DaskalopoulosWentworth}.  Its energy defines a
differentiable function $E_\rho$ on Teichm\"uller space.  The first-variation
formula \cite[Theorem~3.11]{DaskalopoulosWentworth} says that $J$ is critical
precisely when the Hopf differential of $f_{\rho,J}$ vanishes, equivalently
when $f_{\rho,J}$ is a possibly branched minimal map.

We call a representation in $\PU(2,1)$ \emph{irreducible} if its image
preserves no proper complex projective subspace of $\mathbb{CP}^2$.
Equivalently, its image fixes no point of $\mathbb{CP}^2$.  According to the
sign of a representative line for such a point, a reducible subgroup fixes a
point of $\CH^2$, fixes a point of $\partial\CH^2$, or preserves a complex
geodesic (by passing to the orthogonal complement in the last case).

We construct irreducible representations with nonintegral Toledo invariant
whose energy functions have no critical points.  We prove the followoing: 

\begin{theorem}\label{thm:main}
Fix a positive integer $d$ with $3\nmid d$.  There exists $h_0(d)$ such that,
for every $h\ge h_0(d)$ and every $g>h$, there is an irreducible reductive
representation
\[
 \rho_{g,h,d}:\pi_1(S_g)\longrightarrow\PU(2,1)
\]
satisfying
\begin{equation}\label{eq:toledo-main}
 \tau(\rho_{g,h,d})=2h-2-\frac{2d}{3}\notin\mathbb Z,
 \qquad
 \operatorname{Crit}(E_{\rho_{g,h,d}})=\varnothing.
\end{equation}
Branch points are allowed in this assertion.
\end{theorem}

The construction starts with Bronstein's almost-Fuchsian representation
$\rho_0:\pi_1(S_h)\to\PU(2,1)$ and precomposes it with the epimorphism on
fundamental groups induced by an orientation-preserving degree-one pinch
$P:S_g\to S_h$.  Thus $\rho_{g,h,d}=\rho_0\circ P_*$.  The image, and hence
irreducibility, is unchanged, while the new representation has a large kernel.

This construction lies just outside the properness theorem of
Goldman--Wentworth.  They prove that $E_\sigma$ is proper when $\sigma$ is a
convex-cocompact representation \cite[Theorem~A]{GoldmanWentworth}; in their
terminology, this requires $\sigma$ to be an isomorphism of the surface group
onto a convex-cocompact discrete subgroup.  Consequently $E_{\rho_0}$ is
proper on $\Teich(S_h)$ and attains a minimum.  Although
$\rho=\rho_0\circ P_*$ has exactly the same image subgroup as $\rho_0$, it has
nontrivial kernel and is not convex cocompact as a representation of
$\pi_1(S_g)$.  Our conclusion shows that this loss of faithfulness can be
drastic: $E_\rho$ has no critical point at all.  Since a nonnegative proper
energy would attain its infimum, $E_\rho$ cannot be proper.


Nonintegrality alone is not responsible for the conclusion: fractional-Toledo
representations admitting equivariant minimal surfaces are known
\cite{BronsteinPU,LoftinMcIntosh}.  The obstruction here is the combination of
the invariant geometric barrier and the degree-one pinch.  The new point is
that the squared-distance barrier is formulated as a uniform two-plane trace
inequality that applies to weakly conformal harmonic maps even at branch
points.  Equivalently, Theorem~\ref{thm:main} produces points of the
$\PU(2,1)$ character variety outside the image of the forgetful map from
equivariant branched minimal surfaces, an image problem raised in
\cite[Section~6.3]{LoftinMcIntosh}.

\section{Energy, Hopf differentials, and Toledo invariant}

\subsection{The energy function}

Fix an irreducible reductive representation
$\rho:\pi_1(S_g)\to\PU(2,1)$.  For $J\in\Teich(S_g)$, let $f_{\rho,J}$ be
the Corlette harmonic map.  Since we take the irreducible representation, the displacement function of the action on
$\CH^2$ is proper, otherwise an escaping sequence on which the displacements
of a finite generating set remain bounded would determine a common fixed
point of $\im\rho$ on $\partial\CH^2$, contrary to irreducibility.  Thus the
equivariant results cited below apply even when $\rho$ has a kernel.  We use
the convention
\[
 E_\rho(J)=\frac12\int_{S_g}|\dd f_{\rho,J}|^2\,\dd A_J.
\]
The function is independent of the representative metric in the conformal
class.  The Hopf differential
\[
 q_{\rho,J}=\big\langle (f_{\rho,J})_z,(f_{\rho,J})_z\big\rangle\,\dd z^2
 \in H^0(S_g,K_J^2)
\]
is holomorphic.  The first variation is a nonzero universal multiple of the
real part of the Beltrami--quadratic-differential pairing with $q_{\rho,J}$.
More precisely, existence and uniqueness are supplied by the equivariant
harmonic-map theorems, and differentiability and the variation formula are
given in \cite[Theorems~3.2, 3.4, and~3.11]{DaskalopoulosWentworth}.
Since the Beltrami--quadratic-differential pairing is nondegenerate,
\begin{equation}\label{eq:critical-equivalence}
 J\in\operatorname{Crit}(E_\rho)
 \quad\Longleftrightarrow\quad q_{\rho,J}=0
 \quad\Longleftrightarrow\quad f_{\rho,J}
 \text{ is weakly conformal}.
\end{equation}
Since $\rho$ is irreducible, the harmonic map is nonconstant.  A nonconstant
weakly conformal harmonic map from a surface is a branched minimal immersion;
its branch set is discrete
\cite[Propositions~2.2 and~2.4]{GulliverOssermanRoyden}.

\subsection{Normalization and sign of the Toledo invariant}

All complex structures on $S_g$ are compatible with its fixed orientation.
Let $J_{\CH^2}$ and $g_{\CH^2}$ be the complex structure and metric on
$\CH^2$, and take the positive K\"ahler form
\[
 \omega(u,v)=g_{\CH^2}(J_{\CH^2}u,v).
\]
For any smooth $\rho$-equivariant map
$f:\widetilde S_g\to\CH^2$, the invariant form $f^*\omega$ descends to
$S_g$, and we set
\begin{equation}\label{eq:toledo-def}
 \tau(\rho)=\frac{2}{\pi}\int_{S_g}f^*\omega.
\end{equation}
The number in \eqref{eq:toledo-def} is independent of $f$,
and
\begin{equation}\label{eq:toledo-lattice}
 \tau(\rho)\in\frac23\mathbb Z,
 \qquad |\tau(\rho)|\le 2g-2.
\end{equation}
The values index the connected components of the $\PU(2,1)$ character
variety \cite{GKL,Toledo1989,Xia}.  Under this normalization, a representation
lifts to $\SU(2,1)$ precisely when its Toledo invariant belongs to
$2\mathbb Z$ \cite[Remark~2.2]{LoftinMcIntosh}.  Thus the phrase
\emph{nonintegral Toledo invariant} means
\[
 \tau(\rho)=\frac{2k}{3}\quad\text{with}\quad 3\nmid k.
\]

If $\kappa$ is an antiholomorphic isometry of $\CH^2$ and
$\rho^\kappa(\gamma)=\kappa\rho(\gamma)\kappa^{-1}$, then
\begin{equation}\label{eq:toledo-conjugation}
 \tau(\rho^\kappa)=-\tau(\rho),
\end{equation}
because $\kappa^*\omega=-\omega$.

Degree naturality follows directly from the characteristic-class definition.
It is also immediate by pulling back the form in \eqref{eq:toledo-def}.

\begin{lemma}\label{lem:naturality}
Let $P:S_g\to S_h$ be a continuous map of signed degree $m$ between the
oriented surfaces, and let
$\rho_0:\pi_1(S_h)\to\PU(2,1)$.  Then
\[
 \tau(\rho_0\circ P_*)=m\,\tau(\rho_0).
\]
\end{lemma}

\begin{proof}
Replace $P$ by a smooth representative.  If
$F:\widetilde S_h\to\CH^2$ is $\rho_0$-equivariant, then $F^*\omega$ is
$\pi_1(S_h)$-invariant and descends to a two-form $\alpha_0$ on $S_h$.
Moreover, $F\circ\widetilde P$ is $(\rho_0\circ P_*)$-equivariant, and its
pulled-back K\"ahler form descends to $P^*\alpha_0$.  Hence
\[
 \int_{S_g}P^*\alpha_0=m\int_{S_h}\alpha_0.
\]
Multiplying by $2/\pi$ proves the formula.
\end{proof}

\section{The almost-Fuchsian input and the representation}

\subsection{Bronstein's fractional-Toledo representations}

Equip $\CH^2$ with its standard K\"ahler metric $g_{\CH^2}$, normalized so
that its sectional curvatures lie in $[-4,-1]$.  If $J_{\CH^2}$ is its complex
structure, then
\[
 \omega(u,v)=g_{\CH^2}(J_{\CH^2}u,v)
\]
is the positive K\"ahler form used in \eqref{eq:toledo-def}.  All norms and
distances in this subsection are taken with respect to $g_{\CH^2}$ and the
metrics it induces.

For an immersed surface $F:\Sigma\to\CH^2$ and a unit normal vector $\nu$,
let $A_\nu$ be the shape operator defined by
\[
 \langle A_\nu u,v\rangle=\langle\II_F(u,v),\nu\rangle,
\]
and set
\begin{equation}\label{eq:II-operator-norm}
 \|\II_F\|_{\mathrm{op}}
 =\sup_{|\nu|=1}\|A_\nu\|_{\mathrm{op}}
 =\sup_{|u|=|v|=1}|\II_F(u,v)|.
\end{equation}
The full tensor norm dominates this operator norm.  We call $F$
\emph{almost-Fuchsian} if its induced metric is complete and
$\sup_\Sigma\|\II_F\|_{\mathrm{op}}<1$.

In \cite[Section~3]{BronsteinPU}, Bronstein uses the same normalization of the complex-hyperbolic metric as we do: its sectional curvatures lie in
$[-4,-1]$.  His Theorem~6.1 gives an equivariant holomorphic immersion whose
second fundamental form is arbitrarily small.  The pointwise norm used there
and the operator norm defined in \eqref{eq:II-operator-norm} are uniformly equivalent,
since the tangent and normal ranks are fixed; in particular, after decreasing
Bronstein's input parameter if necessary, his estimate implies any prescribed
bound on $\|\II_F\|_{\mathrm{op}}$.  The embedding and convex-cocompactness
conclusions then follow from
\cite[Theorem~3.2 and Corollary~3.3]{BronsteinPU}.  Since a holomorphic curve
in a K\"ahler manifold is minimal, these results give the following form of
Bronstein's construction.

\begin{theorem}[Bronstein, in the present normalization]\label{thm:bronstein}
Fix a positive integer $d$ and $0<\eta<\tfrac12$.  There is
$h_0=h_0(d,\eta)$ such that, for every $h\ge h_0$, there are a complex
structure $J_0$ on $S_h$, a convex-cocompact representation
\[
 \rho_0:\pi_1(S_h)\longrightarrow\PU(2,1),
\]
and a $\rho_0$-equivariant holomorphic embedding
\(
 F:(\widetilde S_h,J_0)\longrightarrow\CH^2
\)
with complete induced metric and
\(
 \sup_{\widetilde S_h}\|\II_F\|_{\mathrm{op}}\le\eta.
\)
The immersion $F$ is minimal and almost-Fuchsian, and
\begin{equation}\label{eq:bronstein-toledo}
 \tau(\rho_0)=2h-2-\frac{2d}{3}.
\end{equation}
\end{theorem}

We quickly explain explicitly the Toledo sign convention.  This is useful
because the convention for the Toledo invariant in \cite{BronsteinPU} is
opposite to the one used here. The discussion in next three paragraphs is just an overview; the notations and terminologies are as in \cite{BronsteinPU}. 

In \cite[Theorem~5.2]{BronsteinPU}, for $h$ sufficiently large,  with given condition, we get a Riemann surface
$(S_h,J_0)$, a holomorphic line bundle $N_h\to (S_h,J_0)$ of degree
\[
    \deg N_h=3h-3+d,
\]
and a nonzero holomorphic section
\(\alpha_h\in H^0(S_h,N_h).
\) Let $K_{J_0}$ denote the canonical bundle of $(S_h,J_0)$.  Bronstein then
introduces the line bundle
\[
    L_h=K_{J_0}^{3}N_h^{-1}.
\]
Note that 
\(K_{J_0}^{3}L_h^{-1}
    =K_{J_0}^{3}\bigl(K_{J_0}^{3}N_h^{-1}\bigr)^{-1}
    \cong N_h.
\)
Thus the section $\alpha_h$ may equivalently be regarded as a section
\(\beta_h\in H^0\bigl(S_h,K_{J_0}^{3}L_h^{-1}\bigr),
\)
which is precisely the holomorphic datum appearing in
\cite[Theorem~3.7]{BronsteinPU}; geometrically, it determines the
$(2,0)$-part of the second fundamental form of the resulting holomorphic
immersion.

\smallskip 

These data determine a
$\rho_0$-equivariant holomorphic immersion
\[
    F:(\widetilde S_h,J_0)\longrightarrow\CH^2
\]
for a representation
\(
    \rho_0:\pi_1(S_h)\longrightarrow\PU(2,1).
\)  In Bronstein's convention,
\cite[equation~(3.5)]{BronsteinPU} gives
\[
    \operatorname{Tol}_{\mathrm B}(\rho_0)
    =-\frac23\deg L_h.
\]
Using the degree computed above,
\[
    \operatorname{Tol}_{\mathrm B}(\rho_0)
    =-\frac23(3h-3-d)
    =2-2h+\frac{2d}{3}.
\]

To relate this formula to our convention, recall the bundle underlying
Bronstein's projectively flat construction.  As a smooth bundle it has the
form
\[
    E_h=K_{J_0}^{-1}\oplus K_{J_0}L_h^{-1}\oplus\mathcal O,
\]
with the Hermitian form of signature $(2,1)$ positive on the first two
summands and negative on the last.  With the holomorphic structure determined
by the solution of the curvature equations, the positive rank-two summand
$V_h$ is represented by the extension
\begin{equation}\label{eq:bronstein-extension}
    0\longrightarrow K_{J_0}^{-1}
    \longrightarrow V_h
    \longrightarrow K_{J_0}L_h^{-1}
    \longrightarrow0.
\end{equation}
We have 
\( \deg V_h =\deg K_{J_0}^{-1}
      +\deg(K_{J_0}L_h^{-1}) =-\deg L_h.
\)

Our normalization of the Toledo invariant is
\(\tau(\rho)
    =\frac{2}{\pi}\int_{S_h}F^*\omega.
\)  This is the convention used in
\cite[Section~2.2, especially equation~(2.9) and the formula preceding
Remark~2.2]{LoftinMcIntosh}.  In this normalization the Chern--Weil formula
for the positive rank-two summand is
\[
    \tau(\rho_0)=-\frac23\deg V_h.
\]
Using $\deg V_h=-\deg L_h$, we therefore obtain
\(\tau(\rho_0) =2h-2-\frac{2d}{3}.
\)
Consequently,
\[
    \tau(\rho_0)
    =-\operatorname{Tol}_{\mathrm B}(\rho_0).
\]


The sign reversal also has a direct geometric explanation.  The map $F$ is
holomorphic and our K\"ahler form is positive.  Hence, away from zero points
since $F$ is an immersion,
\(
    F^*\omega
\)
is the positive area form of the induced metric on $(S_h,J_0)$.  Thus
\( \int_{S_h}F^*\omega>0,\)
so our Toledo invariant must be positive.  This agrees with
\(\tau(\rho_0)=2h-2-\frac{2d}{3}>0
\)
for the range of $h$ used in the construction. 


\bigskip 

For the proof of Theorem~\ref{thm:main}, fix $\eta_0=\tfrac14$ and set
$h_0(d)=h_0(d,\eta_0)$, increasing it if necessary so that
$3h-3-d>0$ for $h\ge h_0(d)$.  Apply Theorem~\ref{thm:bronstein} with
$\eta=\eta_0$ and put
\[
 Y=F(\widetilde S_h)\subset\CH^2.
\]
The surface $Y$ is complete, embedded, minimal, and almost-Fuchsian.  Let
$NY$ be its normal bundle and define
$\exp^\perp(y,\xi)=\exp_y(\xi)$.  Bronstein's normal-exponential theorem
\cite[Theorem~3.1]{BronsteinHadamard} gives a diffeomorphism
\begin{equation}\label{eq:normalexp}
 \exp^\perp:NY\longrightarrow\CH^2.
\end{equation}
In these normal coordinates, the squared distance to $Y$ is
\begin{equation}\label{eq:psi}
 \Psi\bigl(\exp^\perp(y,\xi)\bigr)=|\xi|^2.
\end{equation}
Consequently, $\Psi(x)=d(x,Y)^2$ is smooth on all of $\CH^2$, including along
$Y$.

\subsection{The degree-one pinch}

Fix $g>h$.  There is an orientation-preserving degree-one pinch map
\[
 P:S_g\longrightarrow S_h
\]
that collapses the last $g-h$ handles.  On the standard generators it induces
the epimorphism
\begin{equation}\label{eq:pinch}
 P_*(a_i)=a_i,\quad P_*(b_i)=b_i\quad(i\le h),
 \qquad
 P_*(a_i)=P_*(b_i)=1\quad(i>h).
\end{equation}
Define
\begin{equation}\label{eq:rho-pullback}
 \rho=\rho_0\circ P_*:\pi_1(S_g)\longrightarrow\PU(2,1).
\end{equation}
This is the representation which is going to be the required one as in Theorem \ref{thm:main}.

\begin{proposition}\label{prop:rho-properties}
The representation $\rho$ in \eqref{eq:rho-pullback} is reductive, irreducible, and nonfaithful. Further it satisfies
\[
 \tau(\rho)=2h-2-\frac{2d}{3}.
\]
When $3\nmid d$, its Toledo invariant is nonintegral.
\end{proposition}

\begin{proof}
Since $P_*$ is onto,
\[
 \im\rho=\im\rho_0.
\]
The representation $\rho_0$ is convex-cocompact, hence discrete, faithful,
nonelementary, and reductive.  We prove its irreducibility directly.  If its
image fixed a point of $\CH^2$, it would be a discrete subgroup of a compact
stabilizer and hence finite, contradicting faithfulness.  A fixed point of
$\partial\CH^2$ is likewise incompatible with nonelementary convex
cocompactness.

It remains to rule out preservation of a complex geodesic $C$.  The kernel of
the restriction homomorphism
\[
 \operatorname{Stab}_{\PU(2,1)}(C)
 \longrightarrow\operatorname{Isom}^+(C)
\]
is compact.  Its intersection with the discrete group $\im\rho_0$ is finite
and is trivial because $\im\rho_0\cong\pi_1(S_h)$ is torsion-free.  The
restriction homomorphism is proper because its kernel is compact, so the
induced action on $C$ is consequently discrete and faithful.  It is also
convex-cocompact: an orbit in the totally geodesic $C$ has the same ambient
distances as it has in $C$.  A convex-cocompact Fuchsian group isomorphic to
a closed surface group is cocompact (otherwise its compact convex core has
nonempty boundary and free fundamental group).  Its Toledo invariant must
therefore have absolute value $2h-2$ \cite{Toledo1989}.  This contradicts
\[
 0<\tau(\rho_0)=2h-2-\frac{2d}{3}<2h-2.
\]
The alternatives listed in the definition of irreducibility are exhausted,
so $\rho_0$ is irreducible.  Reductivity and irreducibility depend only on
the image and therefore also hold for $\rho$.

On the other hand, $P_*$ kills the fundamental groups carried by the collapsed
handles, so $\rho$ is not faithful.  Degree naturality
(Lemma~\ref{lem:naturality}) and $\deg P=1$ give the Toledo formula.  Finally,
\[
 \tau(\rho)=\frac23(3h-3-d),
\]
which is not an integer when $3\nmid d$.
\end{proof}

\begin{remark}
The representation $\rho$ is not convex-cocompact as a representation of
$\pi_1(S_g)$: its orbit map cannot be a quasi-isometric embedding because its
kernel is nontrivial.  Its \emph{image subgroup}, however, is the same
convex-cocompact subgroup as that of $\rho_0$.  This distinction is precisely
what makes the construction work.
\end{remark}

\section{A quantitative trapping lemma, including branches}

Bronstein \cite[Proposition~4.5]{BronsteinHadamard} already establishes a
strong squared-distance subharmonicity inequality on an immersed minimal
submanifold in its setting.  Here we give a self-contained pointwise
reformulation tailored to the present application.  It is uniform over every
two-plane, and the min--max
step below explicitly includes planes mixing the radial and tube directions.
This form can be pulled back by a weakly conformal harmonic map without first
assuming that its image is an immersed submanifold.  

\begin{lemma}\label{lem:tube}
Let $X$ be a Hadamard manifold with $K_X\le-1$, and let
$Y^2\subset X$ be a complete embedded minimal surface such that
\[
 \sup_{y\in Y}\|\II_Y\|_{\mathrm{op}}\le\eta<\tfrac12,
 \qquad
 \exp^\perp:NY\longrightarrow X\ \text{is a diffeomorphism}.
\]
Put $\Psi=d(\,\cdot\,,Y)^2$.  Then $\Psi\in C^\infty(X)$ and, for every
$R>0$, there is $c_R>0$ such that, whenever $d(x,Y)\le R$ and
$L\subset T_xX$ is a two-plane,
\begin{equation}\label{eq:tube-estimate}
 \tr_L\Hess\Psi\ge c_R\Psi(x).
\end{equation}
Here $\tr_L$ denotes the trace of the restriction to $L$.
\end{lemma}

\begin{proof}
Under the inverse normal-exponential map, $\Psi(y,v)=|v|^2$, so $\Psi$ is
smooth.  Fix $x\notin Y$, write $r=d(x,Y)$ and $n=\dim X$, and let $\nu$ be
the unit normal at the foot point that points toward $x$.  Let
$\mu_1\le\cdots\le\mu_{n-1}$ be the principal curvatures, with normal
$\nabla r$, of the level hypersurface $N_rY$.  Since $Y$ is minimal, the two
eigenvalues of its shape operator $A_\nu$ are $-a,a$ for some
$0\le a\le\eta$.

Put $t=\tanh r$ and $q_r(s)=(t+s)/(1+st)$.  Bronstein's comparison
\cite[Corollary~3.7]{BronsteinHadamard} gives
\begin{align}
 \mu_1&\ge q_r(-a)\ge q_r(-\eta),\label{eq:mu-one}\\
 \mu_1+\mu_2
 &\ge q_r(-a)+q_r(a)
 =\frac{2t(1-a^2)}{1-a^2t^2}
 \ge 2(1-\eta^2)t.\label{eq:mu-two}
\end{align}
Off $Y$ one has
\begin{equation}\label{eq:hessian-distance}
 \Hess(r^2)=2\,\dd r\otimes\dd r+2r\,\Hess r,
\end{equation}
so its eigenvalues are $2,2r\mu_1,\ldots,2r\mu_{n-1}$.  By the min--max
principle, the trace on any two-plane is at least the sum of the two least
eigenvalues.  It is therefore bounded below by
\begin{equation}\label{eq:mixed-min}
 \min\left\{2r(\mu_1+\mu_2),\ 2+2r\mu_1\right\}.
\end{equation}
This is the step that also treats mixed radial--tube planes.

By \eqref{eq:mu-two}, the first term in \eqref{eq:mixed-min} is at least
\[
 4(1-\eta^2)r\tanh r
 \ge 4(1-\eta^2)\frac{\tanh R}{R}r^2.
\]
Here we used that $r\mapsto\tanh r/r$ is decreasing on $(0,\infty)$.
Moreover, $q_r(-\eta)\ge-\eta$ and it is nonnegative for
$r\ge\operatorname{arctanh}\eta$.  Hence \eqref{eq:mu-one} implies
\[
 2+2r\mu_1\ge2\delta_\eta,
 \qquad
 \delta_\eta=1-\eta\operatorname{arctanh}\eta>0;
\]
in the only nontrivial range $0<r\le\operatorname{arctanh}\eta$, use
$\mu_1\ge-\eta$.  Since $r\le R$, one may take
\begin{equation}\label{eq:cR}
 c_R=\min\left\{
 4(1-\eta^2)\frac{\tanh R}{R},
 \frac{2\delta_\eta}{R^2}
 \right\}>0.
\end{equation}
At a point $x\in Y$, $\Hess\Psi$ is twice the orthogonal projection onto
$N_xY$, while the right-hand side of \eqref{eq:tube-estimate} vanishes.  This
completes the proof.
\end{proof}

The next statement is the form needed for the energy problem.  Notice that the
domain quotient is compact, so no stochastic-completeness hypothesis is
required.  We use the sign convention
\[
 \Delta=\tr\Hess;
\]
in particular, $\Delta u\le0$ at a local maximum of a smooth function.

\begin{lemma}[Branched trapping]\label{lem:trapping}
Assume that $X$ and $Y$ satisfy the hypotheses of Lemma~\ref{lem:tube}.  Let $M$
be a closed connected Riemann surface, let
$\sigma:\pi_1(M)\to\operatorname{Isom}(X)$ have image preserving $Y$, and let
\[
 f:\widetilde M\longrightarrow X
\]
be a nonconstant $\sigma$-equivariant harmonic map which is weakly conformal.
Then $f(\widetilde M)\subset Y$.
\end{lemma}

\begin{proof}
Because $Y$ is $\sigma(\pi_1M)$-invariant, the function
\[
 u=\Psi\circ f=d(f(\,\cdot\,),Y)^2
\]
is $\pi_1(M)$-invariant and descends to a smooth nonnegative function on the
compact surface $M$.  Set
\[
 R_0=\max\left\{1,\max_M\sqrt{u}\right\},
 \qquad e(f)=\frac12|\dd f|^2.
\]

At a point where $e(f)>0$, choose an orthonormal frame $e_1,e_2$ on $M$.
Weak conformality gives
\[
 \dd f(e_i)=\lambda E_i,
\]
where $E_1,E_2$ are orthonormal in a two-plane $L\subset TX$ and
$\lambda^2=e(f)$.  Since $f$ is harmonic, the composition formula and
Lemma~\ref{lem:tube} give
\begin{align}
 \Delta_Mu
 &=\sum_{i=1}^{2}\Hess\Psi(\dd f(e_i),\dd f(e_i))
   +\dd\Psi(\tau(f))\notag\\
 &=\lambda^2\tr_L\Hess\Psi
 \ge c_{R_0}e(f)u.\label{eq:subharmonic}
\end{align}
If $\dd f=0$, the same composition formula gives $\Delta_Mu=0$, so
inequality~\eqref{eq:subharmonic} holds pointwise there as well.  Thus
$\Delta_Mu\ge0$ on the closed connected surface $M$, and the maximum principle
makes $u$ constant.  If this constant were positive,
inequality~\eqref{eq:subharmonic} would force
$e(f)\equiv0$, contrary to the assumption that $f$ is nonconstant.  Hence
$u\equiv0$, which is the assertion.
\end{proof}

\begin{remark}
The use of $r^2$, rather than $r$, is convenient at points of $Y$: it is smooth
under the normal-exponential diffeomorphism.  At a branch point the composition
formula gives $\Delta_Mu=0$ directly, so no limiting argument or choice of a
two-plane is needed.
\end{remark}

\section{Proof that the energy has no critical point}

We now prove Theorem~\ref{thm:main}.  Retain the construction of
Proposition~\ref{prop:rho-properties}.

\begin{proof}[Proof of Theorem~\ref{thm:main}]
Suppose, for contradiction, that $J\in\Teich(S_g)$ is a critical point of
$E_\rho$.  By the equivalences in~\eqref{eq:critical-equivalence}, the
Corlette harmonic map
\[
f=f_{\rho,J}:(\widetilde S_g,J)\longrightarrow\CH^2
\]
is weakly conformal and is therefore a possibly branched minimal map.  It is
nonconstant: otherwise $\im\rho$ would fix its value, contrary to
irreducibility.

The group $\im\rho=\im\rho_0$ preserves the almost-Fuchsian plane
$Y=F(\widetilde S_h)$.  Applying Lemma~\ref{lem:trapping} yields
\begin{equation}\label{eq:image-in-Y}
 f(\widetilde S_g)\subset Y.
\end{equation}
Since $F:\widetilde S_h\to Y$ is an embedding, define
\[
 H=F^{-1}\circ f:\widetilde S_g\longrightarrow\widetilde S_h.
\]
Indeed, injectivity of $F$ and equivariance give
\[
 F(H(\gamma x))=f(\gamma x)=\rho_0(P_*(\gamma))f(x)
 =F(P_*(\gamma)H(x)).
\]
Therefore
\begin{equation}\label{eq:H-equivariance}
 H(\gamma x)=P_*(\gamma)H(x),
 \qquad \gamma\in\pi_1(S_g).
\end{equation}
Consequently $H$ descends to a smooth map
\begin{equation}\label{eq:Hbar}
 \overline H:(S_g,J)\longrightarrow(S_h,J_0)
\end{equation}
which induces $P_*$ on fundamental groups after compatible basepoint choices
(and induces it up to an inner automorphism without such choices).

Equip $\widetilde S_h$ with the induced metric $g_Y=F^*g_{\CH^2}$.  For each
$\gamma\in\pi_1(S_h)$, equivariance and invariance of the ambient metric give
\[
 \gamma^*g_Y=(F\circ\gamma)^*g_{\CH^2}
 =(\rho_0(\gamma)\circ F)^*g_{\CH^2}=g_Y.
\]
Thus $g_Y$ descends to a metric on $S_h$.  Because $F$ is holomorphic, the
conformal class of $g_Y$ is $J_0$.  Since $F$ is an isometric immersion for
$g_Y$ and $f=F\circ H$ is weakly conformal, $H$ is weakly conformal.

For a local orthonormal frame $e_1,e_2$ on the domain, the composition
formula for tension fields is
\begin{equation}\label{eq:tension-composition}
 \tau(F\circ H)=\dd F(\tau(H))+
 \sum_{i=1}^2\II_F\bigl(\dd H(e_i),\dd H(e_i)\bigr).
\end{equation}
Where $\dd H\ne0$, choose an orthonormal frame with
$\dd H(e_i)=\lambda E_i$ and $E_1,E_2$ orthonormal.  Then the last term is
$\lambda^2(\II_F(E_1,E_1)+\II_F(E_2,E_2))$ and vanishes because $F$ is
minimal.  Where $\dd H=0$, it vanishes directly.  Since $F\circ H=f$ is
harmonic and $\dd F$ is injective, equation~\eqref{eq:tension-composition}
implies
\[
 \tau(H)=0.
\]
Thus $\overline H$ is a nonconstant harmonic weakly conformal map between
Riemann surfaces.  Its branch set is discrete, its complement is connected,
and the orientation sign is constant there.  The standard local
characterization therefore shows that $\overline H$ is globally holomorphic
or globally antiholomorphic.

The closed surfaces $S_g$ and $S_h$ are aspherical.  Because
$\overline H_*=P_*$, the maps $\overline H$ and $P$ are homotopic and
\[
 \deg\overline H=\deg P=1.
\]
An antiholomorphic nonconstant map has negative signed degree, so
$\overline H$ is holomorphic.  A nonconstant holomorphic map between compact
connected Riemann surfaces is surjective, and the local degrees over any
fiber are positive integers whose sum is the global degree.  Since
$\deg\overline H=1$, every fiber consists of one point of local degree one.
Thus $\overline H$ is unramified and bijective, hence biholomorphic.
Consequently $g=h$, contradicting $g>h$.  Therefore $E_\rho$ has no critical
point.
\end{proof}

\section{About Forgetful map}\label{rem:fractional-warning}
Let
$\widehat{\mathcal M}^{\mathrm{br,irr}}_{g,\tau}$ be the set of triples
$(J,\sigma,f)$ with the following properties: $J$ is an
orientation-compatible complex structure on the marked $S_g$;
$\sigma:\pi_1(S_g)\to\PU(2,1)$ is irreducible and reductive with
$\tau(\sigma)=\tau$; and
\[
 f:(\widetilde S_g,J)\longrightarrow\CH^2
\]
is a nonconstant, $\sigma$-equivariant, harmonic, weakly conformal map.
Branch points are allowed.  We quotient by marked reparametrization and
simultaneous ambient conjugation as follows.  If
$\phi\in\operatorname{Diff}_0^+(S_g)$, $\widetilde\phi$ is a lift, and
$A\in\PU(2,1)$, then
\begin{equation}\label{eq:minimal-moduli-equivalence}
 (J,\sigma,f)\sim
 \bigl(\phi^*J,\ A(\sigma\circ\phi_*)A^{-1},\
 A\circ f\circ\widetilde\phi\bigr).
\end{equation}
Compatible basepoints are understood; changing them only changes the displayed
data by the same equivalence.  Put
\[
 \mathcal M^{\mathrm{br,irr}}_{g,\tau}
 =\widehat{\mathcal M}^{\mathrm{br,irr}}_{g,\tau}/\!\sim.
\]
Thus $J$ remains marked, reparametrizations isotopic to the identity have been
divided out, and the mapping class group has not been divided out.

Let $\cX_\tau^{\mathrm{irr}}(S_g,\PU(2,1))$ be the set of
$\PU(2,1)$-conjugacy classes $[\sigma]$ of irreducible reductive
representations $\sigma:\pi_1(S_g)\to\PU(2,1)$ with
$\tau(\sigma)=\tau$.  This is the marked irreducible character variety in
Toledo component $\tau$.  The forgetful map considered here is
\begin{equation}\label{eq:forgetful-map}
 \Phi_{g,\tau}:\mathcal M^{\mathrm{br,irr}}_{g,\tau}
 \longrightarrow\cX_\tau^{\mathrm{irr}}(S_g,\PU(2,1)),
 \qquad [(J,\sigma,f)]\longmapsto[\sigma].
\end{equation}
It is well defined because $\phi\in\operatorname{Diff}_0^+(S_g)$ induces an
inner automorphism of $\pi_1(S_g)$.  This is the marked version of the
forgetful map discussed in \cite[Section~6.3]{LoftinMcIntosh}.

\begin{corollary}\label{cor:nonsurjective}
For the triples $(g,h,d)$ in Theorem~\ref{thm:main}, the map
$\Phi_{g,2h-2-2d/3}$ is not surjective.  More precisely,
\[
 [\rho_{g,h,d}]\notin\operatorname{im}\Phi_{g,2h-2-2d/3}.
\]
\end{corollary}

\begin{proof}
By uniqueness of the equivariant harmonic map and the equivalences
in~\eqref{eq:critical-equivalence}, a triple above a character $[\sigma]$
exists precisely when $E_\sigma$ has a critical point.  Apply
Theorem~\ref{thm:main}.
\end{proof}





\begin{thebibliography}{Bron26b}

\bibitem[Bron26a]{BronsteinHadamard}
S.~Bronstein,
\emph{On almost-Fuchsian submanifolds of Hadamard spaces and the asymptotic
Plateau problem},
Ann. Inst. Fourier, online first,  2026, 15~pp.,
\url{https://doi.org/10.5802/aif.3767}.

\bibitem[Bron26b]{BronsteinPU}
S.~Bronstein,
\emph{Almost-Fuchsian representations in $\PU(2,1)$},
arXiv:2411.16261v3, 
\url{https://arxiv.org/abs/2411.16261v3}.

\bibitem[Cor88]{Corlette}
K.~Corlette,
\emph{Flat $G$-bundles with canonical metrics},
J. Differential Geom. \textbf{28} (1988), 361--382,
\url{https://doi.org/10.4310/jdg/1214442469}.

\bibitem[DW07]{DaskalopoulosWentworth}
G.~D.~Daskalopoulos and R.~A.~Wentworth,
\emph{Harmonic maps and Teichm\"uller theory},
in \emph{Handbook of Teichm\"uller Theory}, Vol.~I,
IRMA Lect. Math. Theor. Phys. 11, Eur. Math. Soc., Z\"urich, 2007, 33--109,
\url{https://ems.press/books/irma/35/731}.

\bibitem[GKL01]{GKL}
W.~M.~Goldman, M.~Kapovich, and B.~Leeb,
\emph{Complex hyperbolic manifolds homotopy equivalent to a Riemann surface},
Comm. Anal. Geom. \textbf{9} (2001), 61--95,
\url{https://doi.org/10.4310/CAG.2001.v9.n1.a3}.

\bibitem[GW07]{GoldmanWentworth}
W.~M.~Goldman and R.~A.~Wentworth,
\emph{Energy of twisted harmonic maps of Riemann surfaces},
in \emph{In the Tradition of Ahlfors--Bers, IV},
Contemp. Math. \textbf{432}, Amer. Math. Soc., Providence, RI, 2007, 45--61,
\url{https://doi.org/10.1090/conm/432/08299}.

\bibitem[LM19]{LoftinMcIntosh}
J.~Loftin and I.~McIntosh,
\emph{Equivariant minimal surfaces in $\CH^2$ and their Higgs bundles},
Asian J. Math. \textbf{23} (2019), 71--106,
\url{https://doi.org/10.4310/AJM.2019.v23.n1.a5}; current corrected version:
\url{https://arxiv.org/abs/1510.00553v5}.

\bibitem[LM25]{LoftinMcIntoshCorrection}
J.~Loftin and I.~McIntosh,
\emph{A correction to ``Equivariant minimal surfaces in $\CH^2$ and their
Higgs bundles''},
Asian J. Math. \textbf{29} (2025), no.~4, 559--562,
\url{https://doi.org/10.4310/AJM.251216035015}.

\bibitem[Tol89]{Toledo1989}
D.~Toledo,
\emph{Representations of surface groups in complex hyperbolic space},
J. Differential Geom. \textbf{29} (1989), 125--133,
\url{https://doi.org/10.4310/jdg/1214442638}.

\bibitem[Xia00]{Xia}
E.~Z.~Xia,
\emph{The moduli of flat $\PU(2,1)$ structures on Riemann surfaces},
Pacific J. Math. \textbf{195} (2000), 231--256,
\url{https://doi.org/10.2140/pjm.2000.195.231}.

\end{thebibliography}

\end{document}
